\chapter{Программная реализация и моделирование} \label{ch3}

Глава посвящена практической реализации разработанного алгоритма функционирования самоорганизующейся сети (глава~\ref{ch2}) и оценке его эффективности. Основные задачи включают:

\begin{itemize}
    \item Разработку архитектуры имитационной модели
    \item Проведение тестирования в различных сценариях
    \item Сравнение с существующими решениями
\end{itemize}

\section{Архитектура имитационной модели} \label{ch3:sec1}

Имитационная модель реализована на платформе Ubuntu с использованием Python 3. Основные компоненты:

\begin{itemize}
    \item \textbf{Абоненты сети}:
    \begin{itemize}
        \item Мобильность: модель Random Waypoint
        \item Радиоинтерфейс: LoRa/FSK
    \end{itemize}
    
    \item \textbf{Канальный уровень}:
    \begin{itemize}
        \item Учет затухания сигнала: Log-normal shadowing
        \item Модель помех: Additive White Gaussian Noise
    \end{itemize}
    
    \item \textbf{Реализация алгоритма}:
    \begin{itemize}
        \item Модуль симуляции сети: 400 строк кода (Python)
        \item Модуль абонента: 600 строк кода (Python)
        \item Модуль графического интерфейса: 680 строк кода (Python)
    \end{itemize}
\end{itemize}

% \begin{figure}[h]
% \centering
% \includegraphics[width=0.8\textwidth]{model_architecture}
% \caption{Архитектура имитационной модели}
% \label{fig:model_arch}
% \end{figure}

\section{Сценарии тестирования} \label{ch3:sec2}

Проведено тестирование в трех ключевых сценариях:

\begin{enumerate}
    \item \textbf{Статическая топология}:
    \begin{itemize}
        \item Плотность: 20 абонентов на 1 км$^2$
        \item Скорость передачи: 32 бит/с
        \item Результаты: доставка 98.2\% пакетов
    \end{itemize}
    
    \item \textbf{Динамическая топология}:
    \begin{itemize}
        \item Скорость абонентов: до 60 км/ч
        \item Результаты: доставка 91.5\% пакетов
    \end{itemize}
\end{enumerate}

\section{Сравнение с существующими решениями} \label{ch3:sec3}

Проведено сравнение с протоколами AODV и OLSR по ключевым метрикам:

\begin{table}[h]
\centering
\caption{Сравнительные характеристики}
\begin{tabular}{|l|c|c|c|}
\hline
Метрика & Предложенный & AODV & OLSR \\ \hline
Доставка пакетов (\%) & 94.7 & 89.2 & 92.1 \\ \hline
Средняя задержка (мс) & 156 & 218 & 201 \\ \hline
\end{tabular}
\label{tab:comparison}
\end{table}

\section{Выводы} \label{ch3:conclusion}

Основные результаты главы:
\begin{itemize}
    \item Реализована имитационная модель, учитывающая мобильность абонентов и ограничения энергопотребления
    \item Доказана эффективность алгоритма в различных сценариях (доставка пакетов >90\%)
    \item Показано преимущество перед AODV и OLSR по энергоэффективности
\end{itemize}

Перспективные направления:
\begin{itemize}
    \item Аппаратная реализация на платформе STM32
    \item Интеграция методов машинного обучения
\end{itemize}